\section{黄金存在态：投影、条件化与有效维数}

\subsection{微观态 Hilbert 空间与存在投影}

定义微观态空间 $\Omega_m:=\{0,1\}^m$，并取
\[
\mathcal H_m:=\ell^2(\Omega_m),
\]
以 $\{\lvert w\rangle:w\in\Omega_m\}$ 为标准正交基。定义“黄金存在投影”算子
\[
P_m:=\sum_{w\in X_m}\lvert w\rangle\langle w\rvert.
\]
该算子将任意态投影到“语法合法”的稳定子空间
\[
\mathcal H_m^{\mathrm{stab}}:=P_m\mathcal H_m.
\]

\subsection{黄金存在态的规范化定义}

取均匀叠加态
\[
\lvert\mathrm{unif}\rangle:=2^{-m/2}\sum_{w\in\Omega_m}\lvert w\rangle,
\]
定义黄金存在态为其在稳定子空间上的条件化：
\[
\lvert\Phi_m\rangle:=\frac{P_m\lvert\mathrm{unif}\rangle}{\norm{P_m\lvert\mathrm{unif}\rangle}}
=\frac{1}{\sqrt{\card{X_m}}}\sum_{w\in X_m}\lvert w\rangle.
\]
该定义不依赖主观语言；其含义为：在给定分辨率 $m$ 下，所有合法类型等权出现的规范化参考态。

\subsection{“存在/不存在”作为二值投影读出}

给定任意密度算子 $\rho$（$\mathcal H_m$ 上的状态），定义“存在事件”概率
\[
\mathbb P(\mathrm{exist})=\mathrm{Tr}(P_m\rho),
\]
以及在“存在事件”发生条件下的条件化态
\[
\rho_{\mathrm{obs}}=\frac{P_m\rho P_m}{\mathrm{Tr}(P_m\rho)}.
\]
因此，“只能观测到存在态”可被严格表述为：观测统计以 $\rho_{\mathrm{obs}}$ 为有效态，而非以 $\rho$ 为无条件态。该表述将“存在”从本体断言转化为读出条件化操作。

\subsection{黄金有效维数与熵率}

定义稳定类型增长率（拓扑熵）：
\[
h_{\mathrm{top}}:=\lim_{m\to\infty}\frac{1}{m}\log \card{X_m}.
\]
对黄金均值约束可得 $h_{\mathrm{top}}=\log\varph$，其中 $\varph=\frac{1+\sqrt5}{2}$。由此得到“黄金有效维数”（以二进制单位计）：
\[
d_{\mathrm{eff}}:=\lim_{m\to\infty}\frac{1}{m}\log_2\card{X_m}=\log_2\varph.
\]
该量介于 $0$ 与 $1$ 之间，刻画在 $m$ 比特微观空间中，经语法稳定性筛选后可区分类型的指数尺度。若将“维数”理解为可区分自由度的增长指数，则该指数自然取黄金值。

